\section{Probability as Dependence-Enriched Logic}
\label{sec:probability-dependence-logic}

The previous separation changes how the connection with probability
should be stated.  There is a weak sense in which probability or fuzzy
logic generalizes classical logic: replace the Boolean values
$\{0,1\}$ by values in $[0,1]$.  This is not the sense used here.  A
mere continuum of truth values still leaves open the central question:
what determines the joint value of two propositions?

Fuzzy logics usually answer this by choosing truth functions.  A t-norm
$T(a,b)$ gives the truth value of a conjunction from the two marginal
truth values alone, and a residuum gives a companion implication
\cite{hajek1998,klement2000,cintula2011}.  This is a powerful algebraic
program, but it is not yet probability.  It treats the joint as a fixed
function of the marginals.  In probabilistic language, that is a global
choice of dependence policy.  Product conjunction is the independence
policy, $T(a,b)=ab$.  Minimum conjunction is the comonotone or maximum
positive-dependence policy, $T(a,b)=\min(a,b)$.  Other t-norms encode
other fixed coupling assumptions.  They are not the general joint.

Probability adds exactly the missing degree of freedom.  The marginal
values $P(A)$ and $P(B)$ do not determine $P(A\cap B)$.  The admissible
joint lies in the Fr\'echet interval
\[
  \max(P(A)+P(B)-1,0)\le P(A\cap B)\le \min(P(A),P(B)).
\]
A copula is the corresponding dependence object at distribution level:
Sklar's theorem separates marginal distributions from the dependence
structure that combines them \cite{nelsen2006}.  The event-level moral is
the same: truth values or marginal credences are not enough; the joint is
additional structure.

This gives a sharper path toward the slogan ``probability is extended
logic'' \cite{cox1961,jaynes2003}.  The extension is not simply
\[
  \{0,1\}\leadsto[0,1].
\]
That move gives many-valued semantics.  The probabilistic extension is
better described as
\[
  \text{Boolean truth} \quad+\quad \text{measure} \quad+\quad
  \text{supplied dependence}.
\]
Classical logic is recovered when propositions are crisp and consequence
is evaluated at certainty.  Fuzzy logic is recovered when one fixes
truth-functional connectives on the value scale.  Probability appears
when one refuses to fix the joint from the marginals and instead carries
the dependence data explicitly.

Cred is designed around that refusal.  The product $a\conjc b=ab$ is an
operation on scalar values, not the universal meaning of $A\land B$.
The joint value $j=\cred(A\land B)$ is supplied separately.  The
conditional is then not a residuated truth function but the external
admissibility relation
\[
  \Adm(c;j,e)\quad\Longleftrightarrow\quad ce=j.
\]
At positive evidence, this recovers ordinary conditional probability.
At zero evidence, it exposes the loss of rank instead of selecting a
vacuous truth value.

This also reframes the role of copulas in the present project.  A copula
is not merely an optional statistical add-on.  It is the general answer
to the question that truth-functional fuzzy conjunction answers only in
special cases: how are the marginal values coupled?  The formal results
in the Lean development reflect this boundary.  If one demands a single
truth-functional joint on values satisfying symmetry, idempotence,
Fr\'echet-style bounds, and the copula monotonicity condition, the joint
is forced to be the minimum copula.  Thus a truth-functional and
idempotent ``and'' describes the maximally dependent diagonal world, not
the full probabilistic situation.  Relaxing truth-functionality lets the
joint depend on the propositions themselves, as probability requires.

The resulting hierarchy is:
\[
\begin{array}{rcl}
\text{classical logic} &:& \text{crisp values and crisp consequence},\\[1mm]
\text{many-valued/fuzzy logic} &:& \text{graded values with chosen truth functions},\\[1mm]
\text{Cred} &:& \text{graded values plus external consequence and supplied joints},\\[1mm]
\text{probability} &:& \text{event algebra, measure, conditionalization, and dependence structure}.
\end{array}
\]
The new contribution is the interface between the middle two rows.  It
shows why probability cannot be obtained from fuzzy truth values alone:
one must also free the joint from truth-functionality.  Once that is
done, probability can be viewed as a dependence-enriched generalization
of logic, while Cred supplies the scalar chain-rule kernel of that view.
